% ============================================================================
% ANHANG Q: Rotverschiebung in der IWT
% ============================================================================

\chapter{Rotverschiebung in der IWT}
\label{app:rotverschiebung}

\section{Einleitung}
\label{app:rotverschiebung_einleitung}

Die beobachtete kosmologische Rotverschiebung wird in der WDBT+ \textbf{nicht} als Doppler-Effekt einer Expansion interpretiert, sondern als \textbf{kumulative gravitative Wirkung} auf Photonen während ihrer Reise durch das $\Omega$-Feld (siehe Kapitel~\ref{ch:kosmologie}). 

In der IWT muss dieser Effekt aus den fundamentalen Axiomen abgeleitet werden. 

\textbf{Dieser Anhang zeigt:}
\begin{enumerate}
    \item Die Rotverschiebung in der IWT ist \textbf{keine Dämpfung}.
    \item Die Rotverschiebung in der IWT ist eine \textbf{Massenänderung} der Ladungsträger durch die fraktale Geometrie.
    \item Die Rotverschiebung ist eine \textbf{Energiesenke}, die Energie aus dem Materiesystem entfernt.
    \item Die abgegebene Energie wird an das \textbf{Vakuum} zurückgegeben.
    \item Damit ist der \textbf{geschlossene Energiekreislauf} der IWT vollständig.
\end{enumerate}

\section{Axiomatische Grundlage}
\label{app:rotverschiebung_axiome}

Die Rotverschiebung in der IWT folgt aus \textbf{zwei Axiomen}:

\textbf{Axiom 2 (Informationserhaltung):}
\[
\sum_{k=1}^{N} |I_k|^2 = \text{const} \tag{Q.1}
\]

\textbf{Axiom 6 (Energie als emergenter Informationsfluss):}
Energie ist eine Erhaltungsgröße des diskreten Informationsflusses.

\textbf{Konsequenz:} Die Information bleibt \textbf{erhalten}. Es gibt keine Dämpfung.

\section{Die Masse in der IWT}
\label{app:rotverschiebung_masse}

In der IWT ist die Masse \textbf{keine fundamentale Größe}. Sie ist definiert als:

\[
m_k = \delta \sum_{l \in \mathcal{N}(k)} |I_k - I_l|^2 \tag{Q.2}
\]

(siehe Kapitel~\ref{ch:diskrete_informationsdynamik}, Gleichung 3.8)

Die Masse ist die \textbf{Änderung der Information} zwischen benachbarten Knoten.

\subsection{Masse im fraktalen Raum}

Im fraktalen Raum mit Dimension $D$ skaliert die Masse mit der Distanz $r$:

\[
m(r) = m_0 \cdot \left( \frac{r}{l_0} \right)^{D-3} \tag{Q.3}
\]

\textbf{Begründung:} Die Masse ist ein Maß für die Informationsdichte. Im fraktalen Raum skaliert die Informationsdichte mit $r^{D-3}$.

\section{Die Rotverschiebung am Rand}
\label{app:rotverschiebung_rand}

Das diskrete Netzwerk ist endlich (Axiom 8). Die Ränder werden durch die \textbf{Vakuumenergie des fraktalen Raumes} charakterisiert.

\textbf{Definition:} Die Rotverschiebung am Rand ist die \textbf{Massenänderung} eines Ladungsträgers, der den Rand des Netzwerks erreicht:

\[
z_{\text{Rand}} = \frac{m_{\text{in}}}{m_{\text{out}}} - 1 \tag{Q.4}
\]

Die Amplitude bleibt \textbf{unverändert}:

\[
|I_{\text{out}}|^2 = |I_{\text{in}}|^2 \tag{Q.5}
\]

Die Information bleibt erhalten (Axiom 2).

\section{Die Massenänderung am Rand}
\label{app:rotverschiebung_massenänderung}

Am Rand $r = R$ gilt:

\[
m_{\text{out}} = m_{\text{in}} \cdot \left( \frac{R}{l_0} \right)^{D-3} \tag{Q.6}
\]

Mit $R = L_{Q,0}$ (der Korrelationslänge des Q-Feldes):

\[
m_{\text{out}} = m_{\text{in}} \cdot \left( \frac{L_{Q,0}}{l_0} \right)^{D-3} \tag{Q.7}
\]

\subsection{Die Rotverschiebungsformel}

Aus (Q.4) und (Q.7) folgt:

\[
\boxed{
z = \frac{m_{\text{in}}}{m_{\text{out}}} - 1 = \left( \frac{l_0}{L_{Q,0}} \right)^{D-3} - 1
} \tag{Q.8}
\]

Für $D < 3$ und $l_0 < L_{Q,0}$ ist $z > 0$ – das ist eine \textbf{Rotverschiebung}.

\subsection{Frequenzabhängigkeit}

Die Masse eines Ladungsträgers hängt von seiner Frequenz $\omega$ ab:

\[
m(\omega) = m_0 \cdot \left( \frac{\omega}{\omega_0} \right)^{D-3} \tag{Q.9}
\]

Die Rotverschiebung ist daher frequenzabhängig:

\[
\boxed{
z(\omega) = \left( \frac{\omega}{\omega_{\text{max}}} \right)^{D-3} - 1
} \tag{Q.10}
\]

mit $\omega_{\text{max}} = 2\pi c / l_0$.

\section{Die Rotverschiebung als Energiesenke}
\label{app:rotverschiebung_energiesenke}

Die Energie in der IWT ist (Kapitel 4, Gleichung 4.3):

\[
E = \alpha \cdot |I|^2 + \beta \cdot \left( \frac{d\phi}{dt} \right)^2 \tag{Q.11}
\]

Die Energie setzt sich aus zwei Beiträgen zusammen:
\begin{enumerate}
    \item \textbf{Amplitudenenergie:} $E_A = \alpha \cdot |I|^2$
    \item \textbf{Phasenenergie:} $E_\phi = \beta \cdot (d\phi/dt)^2$
\end{enumerate}

\subsection{Die Kodierung der Masse in der Phase}

Die Masse wird in der Phase kodiert. Damit die Information erhalten bleibt (Axiom 2), muss die Amplitude konstant bleiben. Damit die Energie abnimmt (Energiesenke), muss die Phasenenergie abnehmen.

\textbf{Die Kodierung:}
\[
\phi_{\text{out}} = \phi_{\text{in}} \cdot \frac{m_{\text{in}}}{m_{\text{out}}} \tag{Q.12}
\]

\textbf{Beweis:}
\begin{align*}
\frac{d\phi_{\text{out}}}{dt} &= \frac{m_{\text{in}}}{m_{\text{out}}} \cdot \frac{d\phi_{\text{in}}}{dt} \tag{Q.13} \\
\frac{m_{\text{in}}}{m_{\text{out}}} &< 1 \quad \text{(da } m_{\text{out}} > m_{\text{in}} \text{)} \tag{Q.14} \\
\frac{d\phi_{\text{out}}}{dt} &< \frac{d\phi_{\text{in}}}{dt} \tag{Q.15}
\end{align*}

Die Energie nach der Transformation:
\[
E_{\text{out}} = \alpha \cdot |I_{\text{out}}|^2 + \beta \cdot \left( \frac{d\phi_{\text{out}}}{dt} \right)^2 \tag{Q.16}
\]

Mit (Q.5) und (Q.13):
\[
E_{\text{out}} = \alpha \cdot |I_{\text{in}}|^2 + \beta \cdot \left( \frac{m_{\text{in}}}{m_{\text{out}}} \right)^2 \cdot \left( \frac{d\phi_{\text{in}}}{dt} \right)^2 \tag{Q.17}
\]

Da $\frac{m_{\text{in}}}{m_{\text{out}}} < 1$:
\[
E_{\text{out}} < E_{\text{in}} \tag{Q.18}
\]

\textbf{Die Rotverschiebung ist eine Energiesenke.}

\subsection{Die Energiedifferenz}

Die Energiedifferenz ist:
\[
\Delta E = E_{\text{in}} - E_{\text{out}} = \beta \cdot \left( \frac{d\phi_{\text{in}}}{dt} \right)^2 \cdot \left[ 1 - \left( \frac{m_{\text{in}}}{m_{\text{out}}} \right)^2 \right] > 0 \tag{Q.19}
\]

\section{Die Rückführung zur Vakuumenergie}
\label{app:rotverschiebung_rueckfuehrung}

Die abgegebene Energie wird an das Vakuum zurückgegeben:

\[
E_{\text{vac,out}} = E_{\text{vac,in}} + \Delta E \tag{Q.20}
\]

Das Vakuum wird durch die Fluktuationen repräsentiert (Anhang~\ref{app:unschaerfe_herleitung}). Die Rückführung ist \textbf{explizit} – die Energie wird dem Vakuum zugeführt.

\section{Der geschlossene Energiekreislauf}
\label{app:rotverschiebung_kreislauf}

Die vollständige IWT besteht aus einem geschlossenen Energiekreislauf:

\[
\boxed{
\text{Vakuum} \xrightarrow{\text{Fluktuationen (Anhang P)}} \text{Materie} \xrightarrow{\text{Rotverschiebung (Anhang Q)}} \text{Vakuum}
} \tag{Q.21}
\]

\begin{enumerate}
    \item \textbf{Fluktuationen (Anhang P):} Die diskrete Zeit $T > 0$ erzeugt intrinsische Fluktuationen, die Energie aus dem Vakuum in das Materiesystem transportieren.
    \item \textbf{Strukturbildung:} Die Fluktuationen werden durch die nichtlinearen Terme verstärkt und bilden stabile Muster.
    \item \textbf{Rotverschiebung (Anhang Q):} Am Rand des Netzwerks wird Energie aus dem Materiesystem entfernt und an das Vakuum zurückgegeben.
\end{enumerate}

Die Energiebilanz:
\[
E_{\text{vac}} + E_{\text{Materie}} = \text{const} \tag{Q.22}
\]

Die Information bleibt erhalten (Axiom 2):
\[
\sum |I_k|^2 = \text{const} \tag{Q.23}
\]

Die Energie bleibt erhalten (Axiom 6):
\[
E_{\text{total}} = \text{const} \tag{Q.24}
\]

\section{Die Implementierung in der Simulation}
\label{app:rotverschiebung_implementierung}

Für die Simulation wird die Transformation wie folgt implementiert:

\begin{verbatim}
// Für jeden Randknoten:
double amplitude = sqrt(I_real[i]^2 + I_imag[i]^2);  // bleibt konstant

// Massenänderung am Rand
double m_in = get_mass(I_real[i], I_imag[i]);
double m_out = m_in * pow(l0 / L_Q0, D - 3);

// Phase wird langsamer (Energiesenke)
double phase_in = atan2(I_imag[i], I_real[i]);
double phase_out = phase_in * (m_in / m_out);

// Neue Amplitude (unverändert) und neue Phase
I_real[i] = amplitude * cos(phase_out);
I_imag[i] = amplitude * sin(phase_out);

// Energieverlust ans Vakuum
double E_in = alpha * amplitude^2 + beta * (dphase_in/dt)^2;
double E_out = alpha * amplitude^2 + beta * (dphase_out/dt)^2;
double E_lost = E_in - E_out;
add_to_vacuum(E_lost);

// Informationserhaltung: |I_out|^2 = |I_in|^2
\end{verbatim}

\section{Zusammenfassung}
\label{app:rotverschiebung_zusammenfassung}

\begin{tabular}{l|l}
\hline
\textbf{Aspekt} & \textbf{Formel} \\
\hline
Rotverschiebung (Definition) & $z = \frac{m_{\text{in}}}{m_{\text{out}}} - 1$ \\
Amplitude & $|I_{\text{out}}|^2 = |I_{\text{in}}|^2$ \\
Masse am Rand & $m_{\text{out}} = m_{\text{in}} \cdot \left( \frac{l_0}{L_{Q,0}} \right)^{D-3}$ \\
Phase & $\phi_{\text{out}} = \phi_{\text{in}} \cdot \frac{m_{\text{in}}}{m_{\text{out}}}$ \\
Energie & $E_{\text{out}} < E_{\text{in}}$ \\
Energiedifferenz & $\Delta E = E_{\text{in}} - E_{\text{out}} > 0$ \\
Vakuum & $E_{\text{vac,out}} = E_{\text{vac,in}} + \Delta E$ \\
Informationserhaltung & $\sum |I_k|^2 = \text{const}$ \\
Energieerhaltung & $E_{\text{total}} = \text{const}$ \\
Kreislauf & Vakuum $\rightarrow$ Materie $\rightarrow$ Vakuum \\
\hline
\end{tabular}

\section{Offene Fragen}
\label{app:rotverschiebung_offene_fragen}

\begin{enumerate}
    \item \textbf{Wie wird die Rotverschiebung gemessen?} – Durch den Vergleich der Masse $m$ am Rand und im Zentrum.
    \item \textbf{Wie hängt die Rotverschiebung von der Frequenz ab?} – Sie ist frequenzabhängig: $z(\omega) = (\omega/\omega_{\text{max}})^{D-3} - 1$.
    \item \textbf{Wie hängt die Rotverschiebung mit Photonen zusammen?} – Die IWT kennt keine Photonen. Die beobachtete Rotverschiebung von Licht ist ein emergenter Effekt der Ladungsdynamik.
    \item \textbf{Ist die Energiesenke konsistent mit Axiom 6?} – Ja, weil die Energie nicht vernichtet, sondern ans Vakuum überführt wird.
\end{enumerate}

\textbf{Dieser Anhang ist die theoriekonforme Grundlage für die Implementierung der Rotverschiebung in der IWT.}